















































 et des composantes spécialisées."

ok donc je remplace extension 1 par extension 2 sans problemes?
extension 1
% =================================================
\subsection{Décomposition orthogonale exacte de l'erreur de reconstruction}
% =================================================

Cette sous-section ferme le verrou logique restant : elle montre que la décomposition
\[
\mathcal E_{\mathrm{rec}}(\mathbf P,G)
=
\mathcal E_{\mathrm{coll}}(\mathbf P,G)
+
\mathcal J_{\mathrm{leak}}(\mathbf P)
\]
résulte exactement de la décomposition orthogonale, pour la partition \(\mathbf P\), du résidu
entre l'opérateur cible et son approximation collée bloc-diagonale.

\begin{definition}[Espace des opérateurs et produit de Frobenius]
On pose
\[
\mathfrak E=\mathcal L(\Hh^T),
\]
muni du produit scalaire de Frobenius
\[
\langle A,B\rangle_F=\tr(A^\top B),
\qquad
\|A\|_F^2=\langle A,A\rangle_F.
\]
\end{definition}

\begin{definition}[Projecteurs de blocs]
Pour toute partition orthogonale admissible \(\mathbf P=(P_0,\dots,P_n)\), on définit les projections
bloc-diagonale et hors-diagonale par
\[
\Pi_{\mathrm{coll}}^{\mathbf P}(X)
=
\sum_{i=0}^n P_iXP_i,
\qquad
\Pi_{\mathrm{leak}}^{\mathbf P}(X)
=
\sum_{i\neq j}P_iXP_j,
\qquad
X\in\mathfrak E.
\]
On note aussi
\[
\mathfrak E_{\mathrm{coll}}^{\mathbf P}:=\mathrm{Im}\bigl(\Pi_{\mathrm{coll}}^{\mathbf P}\bigr),
\qquad
\mathfrak E_{\mathrm{leak}}^{\mathbf P}:=\mathrm{Im}\bigl(\Pi_{\mathrm{leak}}^{\mathbf P}\bigr).
\]
\end{definition}

\begin{lemma}[Orthogonalité des blocs]
Les sous-espaces \(\mathfrak E_{ij}^{\mathbf P}:=P_i\mathfrak E P_j\) sont deux à deux orthogonaux pour le produit de Frobenius.
En particulier,
\[
\mathfrak E
=
\bigoplus_{i,j=0}^n \mathfrak E_{ij}^{\mathbf P}
=
\mathfrak E_{\mathrm{coll}}^{\mathbf P}\oplus \mathfrak E_{\mathrm{leak}}^{\mathbf P},
\]
et les projections \(\Pi_{\mathrm{coll}}^{\mathbf P}\) et \(\Pi_{\mathrm{leak}}^{\mathbf P}\) sont orthogonales.
\end{lemma}

\begin{proof}
Soient \(A=P_iXP_j\) et \(B=P_kYP_\ell\). Alors
\[
\langle A,B\rangle_F
=
\tr\!\bigl((P_iXP_j)^\top P_kYP_\ell\bigr)
=
\tr\!\bigl(P_jX^\top P_iP_kYP_\ell\bigr).
\]
Si \(i\neq k\), alors \(P_iP_k=0\), donc \(\langle A,B\rangle_F=0\).

Si \(i=k\) et \(j\neq \ell\), alors
\[
\tr\!\bigl(P_jX^\top P_iYP_\ell\bigr)
=
\tr\!\bigl(P_\ell P_jX^\top P_iY\bigr)=0,
\]
car \(P_\ell P_j=0\). Ainsi, deux blocs distincts sont orthogonaux.

Enfin, comme \(\sum_{i=0}^n P_i=\Id\), tout opérateur \(X\in\mathfrak E\) se décompose en
\[
X=\sum_{i,j=0}^n P_iXP_j,
\]
d'où la somme directe orthogonale annoncée.
\end{proof}

\begin{definition}[Résidu de reconstruction et approximation collée]
On note \(\widehat \Sigma_{\mathrm{coll}}(G)\in\mathfrak E_{\mathrm{coll}}^{\mathbf P}\) l'opérateur reconstruit par la fonction de collage \(G\).

On définit le résidu total de reconstruction par
\[
R_{\mathrm{rec}}(\mathbf P,G)
=
\Sigma-\widehat \Sigma_{\mathrm{coll}}(G).
\]
On définit sa composante collée et sa composante de leakage par
\[
R_{\mathrm{coll}}(\mathbf P,G)
:=
\Pi_{\mathrm{coll}}^{\mathbf P}\!\bigl(R_{\mathrm{rec}}(\mathbf P,G)\bigr),
\qquad
R_{\mathrm{leak}}(\mathbf P)
:=
\Pi_{\mathrm{leak}}^{\mathbf P}(\Sigma).
\]
\end{definition}

\begin{definition}[Erreurs globales]
On définit
\[
\mathcal E_{\mathrm{rec}}(\mathbf P,G)
=
\mathbb E\bigl[\|R_{\mathrm{rec}}(\mathbf P,G)\|_F^2\bigr],
\qquad
\mathcal E_{\mathrm{coll}}(\mathbf P,G)
=
\mathbb E\bigl[\|R_{\mathrm{coll}}(\mathbf P,G)\|_F^2\bigr].
\]
Quand \(\mathbf P\) est fixé, on abrège parfois \(\mathcal E_{\mathrm{coll}}(G)\).
Le terme \(\mathcal E_{\mathrm{coll}}(\mathbf P,G)\) est l'erreur portée par la composante bloc-diagonale reconstruite.
\end{definition}





\begin{definition}[Le terme de leakage]
On pose
\[
\mathcal J_{\mathrm{leak}}(\mathbf P)
=
\bigl\|\Pi_{\mathrm{leak}}^{\mathbf P}(\Sigma)\bigr\|_F^2
=
\sum_{i\neq j}\|P_i\Sigma P_j\|_F^2.
\]
\end{definition}

\begin{theorem}[Décomposition orthogonale exacte de l'erreur de reconstruction]
Sous les définitions précédentes, pour toute architecture admissible \((\mathbf P,G)\),
\[
\mathcal E_{\mathrm{rec}}(\mathbf P,G)
=
\mathcal E_{\mathrm{coll}}(\mathbf P,G)
+
\mathcal J_{\mathrm{leak}}(\mathbf P).
\]
\end{theorem}

\begin{proof}
On écrit d'abord
\[
R_{\mathrm{rec}}(\mathbf P,G)
=
\Sigma-\widehat \Sigma_{\mathrm{coll}}(G).
\]
Comme \(\widehat \Sigma_{\mathrm{coll}}(G)\in \mathfrak E_{\mathrm{coll}}^{\mathbf P}\), sa projection hors-diagonale est nulle, donc
\[
\Pi_{\mathrm{leak}}^{\mathbf P}\!\bigl(R_{\mathrm{rec}}(\mathbf P,G)\bigr)
=
\Pi_{\mathrm{leak}}^{\mathbf P}(\Sigma).
\]
De plus,
\[
R_{\mathrm{rec}}(\mathbf P,G)
=
\underbrace{\Bigl(\Pi_{\mathrm{coll}}^{\mathbf P}(\Sigma)-\widehat \Sigma_{\mathrm{coll}}(G)\Bigr)}_{R_{\mathrm{coll}}(\mathbf P,G)\in\mathfrak E_{\mathrm{coll}}^{\mathbf P}}
+
\underbrace{\Pi_{\mathrm{leak}}^{\mathbf P}(\Sigma)}_{R_{\mathrm{leak}}(\mathbf P)\in\mathfrak E_{\mathrm{leak}}^{\mathbf P}}.
\]
Les deux termes appartiennent à des sous-espaces orthogonaux, donc le terme croisé s'annule exactement :
\[
\|R_{\mathrm{rec}}(\mathbf P,G)\|_F^2
=
\|R_{\mathrm{coll}}(\mathbf P,G)\|_F^2
+
\|R_{\mathrm{leak}}(\mathbf P)\|_F^2.
\]
En prenant l'espérance,
\[
\mathcal E_{\mathrm{rec}}(\mathbf P,G)
=
\mathbb E\bigl[\|R_{\mathrm{coll}}(\mathbf P,G)\|_F^2\bigr]
+
\|R_{\mathrm{leak}}(\mathbf P)\|_F^2.
\]
Par définition,
\[
\mathcal E_{\mathrm{coll}}(\mathbf P,G)
=
\mathbb E\bigl[\|R_{\mathrm{coll}}(\mathbf P,G)\|_F^2\bigr].
\]
En outre,
\[
\|R_{\mathrm{leak}}(\mathbf P)\|_F^2
=
\bigl\|\Pi_{\mathrm{leak}}^{\mathbf P}(\Sigma)\bigr\|_F^2
=
\sum_{i\neq j}\|P_i\Sigma P_j\|_F^2
=
\mathcal J_{\mathrm{leak}}(\mathbf P).
\]
Donc finalement,
\[
\mathcal E_{\mathrm{rec}}(\mathbf P,G)
=
\mathcal E_{\mathrm{coll}}(\mathbf P,G)
+
\mathcal J_{\mathrm{leak}}(\mathbf P),
\]
ce qu'il fallait démontrer.
\end{proof}

\begin{remark}
Le résultat ne repose pas sur une orthogonalité postulée entre deux erreurs abstraites.
Il découle directement de la décomposition orthogonale des opérateurs induite par la partition \(\mathbf P\),
puis de la définition de la reconstruction collée comme approximation bloc-diagonale.
Le terme hors-diagonal de \(\Sigma\) est exactement la part irréductible au collage.
\end{remark}

% =================================================
\subsection{Optimalité conjointe sous compatibilité spectrale exacte}
% =================================================

Cette sous-section établit le résultat structural maximal du formalisme sous une hypothèse de compatibilité exacte entre
(i) le découplage spectral des modules,
(ii) l'ordre hiérarchique des blocs, et
(iii) la reconstruction globale par collage pondéré.

\begin{assumption}[Spectre simple et partitions orthogonales de rangs fixés]
On suppose que la covariance de calibration \(\Sigma\) est symétrique positive semi-définie
et que son spectre est simple.
On ordonne ses valeurs propres de façon décroissante :
\[
\lambda_1>\lambda_2>\cdots>\lambda_{dT}\ge 0,
\]
avec une base orthonormée associée \((e_1,\dots,e_{dT})\) telle que
\[
\Sigma e_k=\lambda_k e_k.
\]

On fixe des tailles entières \(d_0,\dots,d_n\ge 0\) telles que
\[
d_0+\cdots+d_n=dT.
\]
On note
\[
I_0=\{1,\dots,d_0\},\quad
I_1=\{d_0+1,\dots,d_0+d_1\},\quad \dots,\quad
I_n=\{d_0+\cdots+d_{n-1}+1,\dots,dT\}.
\]

On définit la partition canonique
\[
\mathbf P^\star=(P_0^\star,\dots,P_n^\star)
\]
par
\[
P_i^\star=\sum_{k\in I_i} e_ke_k^\top,
\qquad i=0,\dots,n.
\]

On appelle \emph{admissible} toute partition orthogonale
\[
\mathbf P=(P_0,\dots,P_n)
\]
de rangs fixés vérifiant
\[
P_iP_j=0\quad(i\neq j),
\qquad
\sum_{i=0}^n P_i=\Id.
\]
La classe admissible n'est donc pas réduite a priori à \(\{\mathbf P^\star\}\) :
l'unicité sera obtenue par le critère d'optimalité ci-dessous.
\end{assumption}

\begin{assumption}[Compatibilité hiérarchique exacte]
Les architectures admissibles sont les couples \((\mathbf P,G)\) tels que :

\begin{itemize}[leftmargin=2em]
    \item \(\mathbf P=(P_0,\dots,P_n)\) est une partition orthogonale admissible ;
    \item les certificats locaux appartiennent aux sous-espaces associés, c'est-à-dire
    \[
    u_0^{(i)}(w,c)\in \mathrm{Im}(P_0),
    \qquad
    \Psi_i(w,c)\in \mathrm{Im}(P_i);
    \]
    \item le collage global \(G\) est affine en les certificats locaux, et son estimateur partagé
    \(\bar u_0(w,c)\) est non biaisé au sens
    \[
    \mathbb E[\bar u_0(w,c)]=u_0^\star(w,c).
    \]
\end{itemize}
\end{assumption}

\begin{assumption}[Modèle de bruit hétéroscédastique]
Pour chaque \(i\in\{1,\dots,n\}\), les certificats partagés vérifient
\[
u_0^{(i)}=u_0^\star+\varepsilon_i,
\]
où les \(\varepsilon_i\) sont indépendantes, centrées, isotropes, et de variance scalaire
\[
\mathbb E\|\varepsilon_i\|^2=d\,\tau_i,
\qquad \tau_i>0.
\]
\end{assumption}

\begin{definition}[Terme d'ordre spectral]
On fixe des poids hiérarchiques strictement décroissants
\[
\omega_0>\omega_1>\cdots>\omega_n>0.
\]
On définit le terme d'ordre par
\[
\mathcal J_{\mathrm{ord}}(\mathbf P)
=
-\sum_{i=0}^n \omega_i\,\mathrm{tr}(P_i\Sigma).
\]
Ce terme impose que les plus grandes directions spectrales soient affectées aux blocs les plus prioritaires.
\end{definition}

\begin{definition}[Objectif architectural conjoint]
On définit l'objectif global par
\[
\mathcal J_{\mathrm{arch}}(\mathbf P,G)
=
\lambda_{\mathrm{leak}}\,\mathcal J_{\mathrm{leak}}(\mathbf P)
+
\lambda_{\mathrm{ord}}\,\mathcal J_{\mathrm{ord}}(\mathbf P)
+
\lambda_{\mathrm{rec}}\,\mathcal E_{\mathrm{rec}}(\mathbf P,G).
\]
En utilisant le théorème précédent, on obtient l'écriture équivalente
\[
\mathcal J_{\mathrm{arch}}(\mathbf P,G)
=
(\lambda_{\mathrm{leak}}+\lambda_{\mathrm{rec}})\,
\mathcal J_{\mathrm{leak}}(\mathbf P)
+
\lambda_{\mathrm{ord}}\,\mathcal J_{\mathrm{ord}}(\mathbf P)
+
\lambda_{\mathrm{rec}}\,\mathcal E_{\mathrm{coll}}(\mathbf P,G).
\]
\end{definition}

\begin{lemma}[Zéro leakage et compatibilité spectrale]
Pour toute partition admissible \(\mathbf P\),
\[
\mathcal J_{\mathrm{leak}}(\mathbf P)\ge 0.
\]
De plus,
\[
\mathcal J_{\mathrm{leak}}(\mathbf P)=0
\quad\Longleftrightarrow\quad
P_i\Sigma=\Sigma P_i
\ \text{ pour tout } i,
\]
c'est-à-dire si et seulement si chaque bloc \(P_i\) est une somme de projecteurs spectraux de \(\Sigma\).
\end{lemma}

\begin{proof}
Par définition, \(\mathcal J_{\mathrm{leak}}\) est une somme de normes au carré, donc elle est non négative.
Si elle vaut \(0\), alors \(P_i\Sigma P_j=0\) pour tous \(i\neq j\), ce qui signifie que \(\Sigma\) est bloc-diagonale dans la décomposition \(\mathbf P\), donc que chaque \(P_i\) commute avec \(\Sigma\).
Réciproquement, si chaque \(P_i\) commute avec \(\Sigma\), alors les blocs hors diagonale sont nuls et \(\mathcal J_{\mathrm{leak}}(\mathbf P)=0\).
\end{proof}

\begin{lemma}[Réécriture télescopique du terme d'ordre]
Soient
\[
m_r=\sum_{i=0}^r d_i
\qquad\text{et}\qquad
Q_r=\sum_{i=0}^r P_i,
\quad r=0,\dots,n,
\]
avec la convention \(Q_{-1}=0\).
Alors
\[
\sum_{i=0}^n \omega_i\,\mathrm{tr}(P_i\Sigma)
=
\omega_n\,\mathrm{tr}(\Sigma)
+
\sum_{r=0}^{n-1}(\omega_r-\omega_{r+1})\,\mathrm{tr}(Q_r\Sigma).
\]
\end{lemma}

\begin{proof}
Comme \(P_i=Q_i-Q_{i-1}\), on a
\[
\sum_{i=0}^n \omega_i\,\mathrm{tr}(P_i\Sigma)
=
\sum_{i=0}^n \omega_i\bigl(\mathrm{tr}(Q_i\Sigma)-\mathrm{tr}(Q_{i-1}\Sigma)\bigr).
\]
Un regroupement télescopique donne
\[
\sum_{i=0}^n \omega_i\,\mathrm{tr}(P_i\Sigma)
=
\omega_n\,\mathrm{tr}(Q_n\Sigma)
+
\sum_{r=0}^{n-1}(\omega_r-\omega_{r+1})\,\mathrm{tr}(Q_r\Sigma).
\]
Comme \(Q_n=\Id\), on obtient la formule annoncée.
\end{proof}

\begin{lemma}[Sélection canonique par l'ordre spectral]
Parmi les partitions admissibles, le terme
\[
\mathcal J_{\mathrm{ord}}(\mathbf P)
=
-\sum_{i=0}^n \omega_i\,\mathrm{tr}(P_i\Sigma)
\]
est strictement minimisé par la partition canonique \(\mathbf P^\star\), et seulement par elle.
\end{lemma}

\begin{proof}
Par le lemme précédent, minimiser \(\mathcal J_{\mathrm{ord}}\) revient à maximiser chaque quantité \(\mathrm{tr}(Q_r\Sigma)\), puisque les coefficients \(\omega_r-\omega_{r+1}\) sont strictement positifs.

Pour tout projecteur orthogonal \(Q\) de rang \(m_r\), le principe de Ky Fan donne
\[
\mathrm{tr}(Q\Sigma)\le \sum_{k=1}^{m_r}\lambda_k,
\]
avec égalité si et seulement si
\[
Q=\sum_{k=1}^{m_r} e_ke^\top_k,
\]
l'unicité étant assurée par la simplicité du spectre.

En l'appliquant à chaque \(Q_r\), on obtient
\[
\mathrm{tr}(Q_r\Sigma)\le \sum_{k=1}^{m_r}\lambda_k,
\]
et l'égalité simultanée pour tout \(r\) impose
\[
Q_r=\sum_{k=1}^{m_r} e_ke^\top_k
\qquad\text{pour tout } r.
\]
Par différence,
\[
P_i=Q_i-Q_{i-1}
=
\sum_{k\in I_i} e_ke^\top_k
=
P_i^\star.
\]
Donc \(\mathbf P^\star\) est l'unique minimiseur de \(\mathcal J_{\mathrm{ord}}\) sur la classe admissible.
\end{proof}

\begin{proposition}[Optimalité du collage inverse-variance]
Sous le modèle de bruit hétéroscédastique, tout estimateur affine non biaisé de la forme
\[
\tilde u_0=\sum_{i=1}^n \beta_i u_0^{(i)},
\qquad
\sum_{i=1}^n \beta_i=1,
\]
vérifie
\[
\mathbb E\|\tilde u_0-u_0^\star\|^2
=
d\sum_{i=1}^n \beta_i^2\tau_i.
\]
Cette quantité est strictement minimisée par
\[
\beta_i^\star
=
\frac{\tau_i^{-1}}{\sum_{j=1}^n\tau_j^{-1}},
\qquad i=1,\dots,n.
\]
\end{proposition}

\begin{proof}
On a
\[
\mathbb E[\tilde u_0-u_0^\star]
=
\sum_{i=1}^n \beta_i\mathbb E[\varepsilon_i]
=
0.
\]
Par indépendance et isotropie,
\[
\mathbb E\|\tilde u_0-u_0^\star\|^2
=
\mathbb E\left\|\sum_{i=1}^n\beta_i\varepsilon_i\right\|^2
=
\sum_{i=1}^n\beta_i^2\,\mathbb E\|\varepsilon_i\|^2
=
d\sum_{i=1}^n\beta_i^2\tau_i.
\]
La fonction
\[
\beta\mapsto d\sum_{i=1}^n\beta_i^2\tau_i
\]
est strictement convexe sur l'hyperplan \(\sum_i\beta_i=1\).  
Le minimiseur unique est donc obtenu par multiplicateurs de Lagrange, ce qui donne
\[
\beta_i^\star
=
\frac{\tau_i^{-1}}{\sum_{j=1}^n\tau_j^{-1}}.
\]
\end{proof}

\begin{theorem}[Optimalité conjointe et unicité du couple partition spectrale + collage pondéré]
Sous les hypothèses précédentes, le problème
\[
\min_{(\mathbf P,G)\ \mathrm{admissible}} \ \mathcal J_{\mathrm{arch}}(\mathbf P,G)
\]
admet un unique minimiseur admissible.
Ce minimiseur est précisément le couple
\[
(\mathbf P^\star,G^\star),
\]
où :

\begin{itemize}[leftmargin=2em]
    \item \(\mathbf P^\star\) est la partition spectrale canonique ;
    \item \(G^\star\) est le collage pondéré inverse-variance, i.e.
    \[
    \bar u_0^\star(w,c)
    =
    \sum_{i=1}^n \alpha_i^\star(w,c)\,u_0^{(i)}(w,c),
    \qquad
    \alpha_i^\star(w,c)
    =
    \frac{\tau_i^{-1}}{\sum_{j=1}^n \tau_j^{-1}}.
    \]
\end{itemize}

En particulier, le couple \((\mathbf P^\star,G^\star)\) est l'unique architecture admissible qui minimise simultanément le leakage inter-modules et l'erreur globale de reconstruction.
\end{theorem}

\begin{proof}
Posons
\[
A=\lambda_{\mathrm{leak}}+\lambda_{\mathrm{rec}}>0.
\]
Alors, pour tout couple admissible \((\mathbf P,G)\),
\[
\mathcal J_{\mathrm{arch}}(\mathbf P,G)
=
A\,\mathcal J_{\mathrm{leak}}(\mathbf P)
+
\lambda_{\mathrm{ord}}\,\mathcal J_{\mathrm{ord}}(\mathbf P)
+
\lambda_{\mathrm{rec}}\,\mathcal E_{\mathrm{coll}}(\mathbf P,G).
\]

\medskip
\noindent
\textbf{Étape 1 : optimisation de la partie partition.}
Comme \(\mathcal J_{\mathrm{leak}}(\mathbf P)\ge 0\) et que \(\mathbf P^\star\) minimise strictement \(\mathcal J_{\mathrm{ord}}(\mathbf P)\) sur la classe admissible, on a pour toute partition admissible
\[
A\,\mathcal J_{\mathrm{leak}}(\mathbf P)+\lambda_{\mathrm{ord}}\,\mathcal J_{\mathrm{ord}}(\mathbf P)
\ge
\lambda_{\mathrm{ord}}\,\mathcal J_{\mathrm{ord}}(\mathbf P)
\ge
\lambda_{\mathrm{ord}}\,\mathcal J_{\mathrm{ord}}(\mathbf P^\star).
\]
De plus,
\[
\mathcal J_{\mathrm{leak}}(\mathbf P^\star)=0
\]
par compatibilité spectrale de la partition canonique.
Ainsi, la partie dépendant de \(\mathbf P\) est minimisée de manière unique par \(\mathbf P^\star\).

\medskip
\noindent
\textbf{Étape 2 : optimisation du collage.}
Une fois \(\mathbf P^\star\) fixé, la partie dépendant de \(G\) se réduit à \(\mathcal E_{\mathrm{coll}}(\mathbf P^\star,G)\), que l'on abrège en \(\mathcal E_{\mathrm{coll}}(G)\).
Sous le modèle de bruit hétéroscédastique, la proposition précédente montre que l'unique collage affine non biaisé minimisant cette erreur est le collage inverse-variance \(G^\star\), avec
\[
\alpha_i^\star(w,c)=\frac{\tau_i^{-1}}{\sum_{j=1}^n\tau_j^{-1}}.
\]

\medskip
\noindent
\textbf{Étape 3 : unicité globale.}
Si \((\mathbf P,G)\) est un minimiseur global, alors il doit satisfaire simultanément
\[
\mathbf P=\mathbf P^\star
\qquad\text{et}\qquad
G=G^\star,
\]
car l'égalité dans les deux inégalités précédentes impose respectivement l'égalité dans le lemme spectral et l'égalité dans la proposition sur le collage.
Réciproquement, \((\mathbf P^\star,G^\star)\) atteint cette borne minimale.
Le minimiseur global est donc unique et vaut exactement \((\mathbf P^\star,G^\star)\).
\end{proof}

\begin{remark}
Le théorème précédent remplace l'hypothèse de décomposition leakage--reconstruction utilisée dans les versions préliminaires du formalisme.
Le couple \((\mathbf P^\star,G^\star)\) minimise simultanément le terme de découplage spectral et le terme de reconstruction collée dans la classe des architectures admissibles.
\end{remark}

\begin{corollary}[Optimalité conjointe]
Le formalisme PCH-ESM réalise simultanément :
\begin{enumerate}[leftmargin=2em]
    \item le découplage maximal des modules via la partition spectrale canonique ;
    \item la reconstruction globale de risque minimal via le collage pondéré inverse-variance.
\end{enumerate}
Autrement dit, sous les hypothèses précédentes, le couple \((\mathbf P^\star,G^\star)\) est l'unique minimiseur simultané du leakage inter-modules et de l'erreur de reconstruction dans la classe des architectures admissibles.
\end{corollary}






extension 2
% =================================================
\subsection{Décomposition orthogonale exacte de l'erreur de reconstruction}
% =================================================

Cette sous-section ferme le verrou logique restant : elle montre que la décomposition
\[
\mathcal E_{\mathrm{rec}}(\mathbf P,G)
=
\mathcal E_{\mathrm{coll}}(\mathbf P,G)
+
\mathcal J_{\mathrm{leak}}(\mathbf P)
\]
résulte exactement de la décomposition orthogonale, pour la partition \(\mathbf P\), du résidu
entre l'opérateur cible et son approximation collée bloc-diagonale.

\begin{definition}[Espace des opérateurs et produit de Frobenius]
On pose
\[
\mathfrak E=\mathcal L(\Hh^T),
\]
muni du produit scalaire de Frobenius
\[
\langle A,B\rangle_F=\tr(A^\top B),
\qquad
\|A\|_F^2=\langle A,A\rangle_F.
\]
\end{definition}

\begin{definition}[Projecteurs de blocs]
Pour toute partition orthogonale admissible \(\mathbf P=(P_0,\dots,P_n)\), on définit les projections
bloc-diagonale et hors-diagonale par
\[
\Pi_{\mathrm{coll}}^{\mathbf P}(X)
=
\sum_{i=0}^n P_iXP_i,
\qquad
\Pi_{\mathrm{leak}}^{\mathbf P}(X)
=
\sum_{i\neq j}P_iXP_j,
\qquad
X\in\mathfrak E.
\]
On note aussi
\[
\mathfrak E_{\mathrm{coll}}^{\mathbf P}:=\mathrm{Im}\bigl(\Pi_{\mathrm{coll}}^{\mathbf P}\bigr),
\qquad
\mathfrak E_{\mathrm{leak}}^{\mathbf P}:=\mathrm{Im}\bigl(\Pi_{\mathrm{leak}}^{\mathbf P}\bigr).
\]
\end{definition}

\begin{lemma}[Orthogonalité des blocs]
Les sous-espaces \(\mathfrak E_{ij}^{\mathbf P}:=P_i\mathfrak E P_j\) sont deux à deux orthogonaux pour le produit de Frobenius.
En particulier,
\[
\mathfrak E
=
\bigoplus_{i,j=0}^n \mathfrak E_{ij}^{\mathbf P}
=
\mathfrak E_{\mathrm{coll}}^{\mathbf P}\oplus \mathfrak E_{\mathrm{leak}}^{\mathbf P},
\]
et les projections \(\Pi_{\mathrm{coll}}^{\mathbf P}\) et \(\Pi_{\mathrm{leak}}^{\mathbf P}\) sont orthogonales.
\end{lemma}

\begin{proof}
Soient \(A=P_iXP_j\) et \(B=P_kYP_\ell\). Alors
\[
\langle A,B\rangle_F
=
\tr\!\bigl((P_iXP_j)^\top P_kYP_\ell\bigr)
=
\tr\!\bigl(P_jX^\top P_iP_kYP_\ell\bigr).
\]
Si \(i\neq k\), alors \(P_iP_k=0\), donc \(\langle A,B\rangle_F=0\).

Si \(i=k\) et \(j\neq \ell\), alors
\[
\tr\!\bigl(P_jX^\top P_iYP_\ell\bigr)
=
\tr\!\bigl(P_\ell P_jX^\top P_iY\bigr)=0,
\]
car \(P_\ell P_j=0\). Ainsi, deux blocs distincts sont orthogonaux.

Enfin, comme \(\sum_{i=0}^n P_i=\Id\), tout opérateur \(X\in\mathfrak E\) se décompose en
\[
X=\sum_{i,j=0}^n P_iXP_j,
\]
d'où la somme directe orthogonale annoncée.
\end{proof}

\begin{definition}[Résidu de reconstruction et approximation collée]
On note \(\widehat \Sigma_{\mathrm{coll}}(G)\in\mathfrak E_{\mathrm{coll}}^{\mathbf P}\) l'opérateur reconstruit par la fonction de collage \(G\).

On définit le résidu total de reconstruction par
\[
R_{\mathrm{rec}}(\mathbf P,G)
=
\Sigma-\widehat \Sigma_{\mathrm{coll}}(G).
\]
On définit sa composante collée et sa composante de leakage par
\[
R_{\mathrm{coll}}(\mathbf P,G)
:=
\Pi_{\mathrm{coll}}^{\mathbf P}\!\bigl(R_{\mathrm{rec}}(\mathbf P,G)\bigr),
\qquad
R_{\mathrm{leak}}(\mathbf P)
:=
\Pi_{\mathrm{leak}}^{\mathbf P}(\Sigma).
\]
\end{definition}

\begin{definition}[Erreurs globales]
On définit
\[
\mathcal E_{\mathrm{rec}}(\mathbf P,G)
=
\mathbb E\bigl[\|R_{\mathrm{rec}}(\mathbf P,G)\|_F^2\bigr],
\qquad
\mathcal E_{\mathrm{coll}}(\mathbf P,G)
=
\mathbb E\bigl[\|R_{\mathrm{coll}}(\mathbf P,G)\|_F^2\bigr].
\]
Quand \(\mathbf P\) est fixé, on abrège parfois \(\mathcal E_{\mathrm{coll}}(G)\).
Le terme \(\mathcal E_{\mathrm{coll}}(\mathbf P,G)\) est l'erreur portée par la composante bloc-diagonale reconstruite.

L'espérance est prise relativement à la distribution des données de calibration utilisées pour construire
\(\widehat \Sigma_{\mathrm{coll}}(G)\).
\end{definition}

\begin{definition}[Le terme de leakage]
On pose
\[
\mathcal J_{\mathrm{leak}}(\mathbf P)
=
\bigl\|\Pi_{\mathrm{leak}}^{\mathbf P}(\Sigma)\bigr\|_F^2
=
\sum_{i\neq j}\|P_i\Sigma P_j\|_F^2.
\]
Le terme de leakage dépend uniquement de la covariance de calibration \(\Sigma\), supposée fixée ; il est donc déterministe.
\end{definition}

\begin{assumption}[Compatibilité minimale du mécanisme de collage]
Pour toute partition admissible \(\mathbf P\), il existe une constante \(C\ge 0\), indépendante de \(\mathbf P\), et une fonction strictement croissante
\[
\phi_{\mathbf P}:[0,\infty)\to[0,\infty),
\]
telles que, pour tout collage admissible \(G\),
\[
\mathcal E_{\mathrm{coll}}(\mathbf P,G)
=
C
+
\phi_{\mathbf P}\!\left(\mathbb E\bigl[\|\bar u_0(w,c)-u_0^\star(w,c)\|^2\bigr]\right).
\]
\end{assumption}

\begin{proposition}[Principe de réduction du collage]
Sous l'hypothèse précédente, pour toute partition admissible \(\mathbf P\),
\[
\argmin_G \mathcal E_{\mathrm{coll}}(\mathbf P,G)
=
\argmin_G \mathbb E\bigl[\|\bar u_0(w,c)-u_0^\star(w,c)\|^2\bigr].
\]
\end{proposition}

\begin{proof}
Comme \(\phi_{\mathbf P}\) est strictement croissante, minimiser \(\mathcal E_{\mathrm{coll}}(\mathbf P,G)\) équivaut exactement à minimiser
\[
\mathbb E\bigl[\|\bar u_0(w,c)-u_0^\star(w,c)\|^2\bigr].
\]
La constante \(C\) est indépendante de \(G\) et de \(\mathbf P\), donc elle n'affecte pas l'ensemble des minimiseurs.
\end{proof}

\begin{theorem}[Décomposition orthogonale exacte de l'erreur de reconstruction]
Sous les définitions précédentes, pour toute architecture admissible \((\mathbf P,G)\),
\[
\mathcal E_{\mathrm{rec}}(\mathbf P,G)
=
\mathcal E_{\mathrm{coll}}(\mathbf P,G)
+
\mathcal J_{\mathrm{leak}}(\mathbf P).
\]
\end{theorem}

\begin{proof}
On écrit d'abord
\[
R_{\mathrm{rec}}(\mathbf P,G)
=
\Sigma-\widehat \Sigma_{\mathrm{coll}}(G).
\]
Comme \(\widehat \Sigma_{\mathrm{coll}}(G)\in \mathfrak E_{\mathrm{coll}}^{\mathbf P}\), sa projection hors-diagonale est nulle, donc
\[
\Pi_{\mathrm{leak}}^{\mathbf P}\!\bigl(R_{\mathrm{rec}}(\mathbf P,G)\bigr)
=
\Pi_{\mathrm{leak}}^{\mathbf P}(\Sigma).
\]
De plus,
\[
R_{\mathrm{rec}}(\mathbf P,G)
=
\underbrace{\Bigl(\Pi_{\mathrm{coll}}^{\mathbf P}(\Sigma)-\widehat \Sigma_{\mathrm{coll}}(G)\Bigr)}_{R_{\mathrm{coll}}(\mathbf P,G)\in\mathfrak E_{\mathrm{coll}}^{\mathbf P}}
+
\underbrace{\Pi_{\mathrm{leak}}^{\mathbf P}(\Sigma)}_{R_{\mathrm{leak}}(\mathbf P)\in\mathfrak E_{\mathrm{leak}}^{\mathbf P}}.
\]
Les deux termes appartiennent à des sous-espaces orthogonaux, donc le terme croisé s'annule exactement :
\[
\|R_{\mathrm{rec}}(\mathbf P,G)\|_F^2
=
\|R_{\mathrm{coll}}(\mathbf P,G)\|_F^2
+
\|R_{\mathrm{leak}}(\mathbf P)\|_F^2.
\]
En prenant l'espérance,
\[
\mathcal E_{\mathrm{rec}}(\mathbf P,G)
=
\mathbb E\bigl[\|R_{\mathrm{coll}}(\mathbf P,G)\|_F^2\bigr]
+
\|R_{\mathrm{leak}}(\mathbf P)\|_F^2.
\]
Par définition,
\[
\mathcal E_{\mathrm{coll}}(\mathbf P,G)
=
\mathbb E\bigl[\|R_{\mathrm{coll}}(\mathbf P,G)\|_F^2\bigr].
\]
En outre,
\[
\|R_{\mathrm{leak}}(\mathbf P)\|_F^2
=
\bigl\|\Pi_{\mathrm{leak}}^{\mathbf P}(\Sigma)\bigr\|_F^2
=
\sum_{i\neq j}\|P_i\Sigma P_j\|_F^2
=
\mathcal J_{\mathrm{leak}}(\mathbf P).
\]
Donc finalement,
\[
\mathcal E_{\mathrm{rec}}(\mathbf P,G)
=
\mathcal E_{\mathrm{coll}}(\mathbf P,G)
+
\mathcal J_{\mathrm{leak}}(\mathbf P),
\]
ce qu'il fallait démontrer.
\end{proof}

\begin{remark}
La décomposition orthogonale précédente est indépendante du mécanisme de collage.
En revanche, l'identification du minimiseur de \(\mathcal E_{\mathrm{coll}}\) avec le collage inverse-variance repose uniquement sur l'hypothèse minimale de compatibilité du mécanisme de collage, qui exprime que \(\mathcal E_{\mathrm{coll}}\) est une transformation strictement croissante de l'erreur de partage.
Le théorème précédent remplace l'hypothèse de décomposition leakage--reconstruction utilisée dans les versions préliminaires du formalisme.
\end{remark}

% =================================================
\subsection{Optimalité conjointe sous compatibilité spectrale exacte}
% =================================================

Cette sous-section établit le résultat structural maximal du formalisme sous une hypothèse de compatibilité exacte entre
(i) le découplage spectral des modules,
(ii) l'ordre hiérarchique des blocs, et
(iii) la reconstruction globale par collage pondéré.

\begin{assumption}[Spectre simple et partitions orthogonales de rangs fixés]
On suppose que la covariance de calibration \(\Sigma\) est symétrique positive semi-définie
et que son spectre est simple.
On ordonne ses valeurs propres de façon décroissante :
\[
\lambda_1>\lambda_2>\cdots>\lambda_{dT}\ge 0,
\]
avec une base orthonormée associée \((e_1,\dots,e_{dT})\) telle que
\[
\Sigma e_k=\lambda_k e_k.
\]

On fixe des tailles entières \(d_0,\dots,d_n\ge 0\) telles que
\[
d_0+\cdots+d_n=dT.
\]
On note
\[
I_0=\{1,\dots,d_0\},\quad
I_1=\{d_0+1,\dots,d_0+d_1\},\quad \dots,\quad
I_n=\{d_0+\cdots+d_{n-1}+1,\dots,dT\}.
\]

On définit la partition canonique
\[
\mathbf P^\star=(P_0^\star,\dots,P_n^\star)
\]
par
\[
P_i^\star=\sum_{k\in I_i} e_ke^\top_k,
\qquad i=0,\dots,n.
\]

On appelle \emph{admissible} toute partition orthogonale
\[
\mathbf P=(P_0,\dots,P_n)
\]
de rangs fixés vérifiant
\[
P_iP_j=0\quad(i\neq j),
\qquad
\sum_{i=0}^n P_i=\Id.
\]
La classe admissible n'est donc pas réduite a priori à \(\{\mathbf P^\star\}\) :
l'unicité sera obtenue par le critère d'optimalité ci-dessous.
\end{assumption}

\begin{assumption}[Compatibilité hiérarchique exacte]
Les architectures admissibles sont les couples \((\mathbf P,G)\) tels que :

\begin{itemize}[leftmargin=2em]
    \item \(\mathbf P=(P_0,\dots,P_n)\) est une partition orthogonale admissible ;
    \item les certificats locaux appartiennent aux sous-espaces associés, c'est-à-dire
    \[
    u_0^{(i)}(w,c)\in \mathrm{Im}(P_0),
    \qquad
    \Psi_i(w,c)\in \mathrm{Im}(P_i);
    \]
    \item le collage global \(G\) est affine en les certificats locaux, et son estimateur partagé
    \(\bar u_0(w,c)\) est non biaisé au sens
    \[
    \mathbb E[\bar u_0(w,c)]=u_0^\star(w,c).
    \]
\end{itemize}
\end{assumption}

\begin{assumption}[Modèle de bruit hétéroscédastique]
Pour chaque \(i\in\{1,\dots,n\}\), les certificats partagés vérifient
\[
u_0^{(i)}=u_0^\star+\varepsilon_i,
\]
où les \(\varepsilon_i\) sont indépendantes, centrées, isotropes, et de variance scalaire
\[
\mathbb E\|\varepsilon_i\|^2=d\,\tau_i,
\qquad \tau_i>0.
\]
\end{assumption}

\begin{assumption}[Poids strictement positifs]
On suppose
\[
\lambda_{\mathrm{leak}}>0,
\qquad
\lambda_{\mathrm{ord}}>0,
\qquad
\lambda_{\mathrm{rec}}>0.
\]
\end{assumption}

\begin{definition}[Terme d'ordre spectral]
On fixe des poids hiérarchiques strictement décroissants
\[
\omega_0>\omega_1>\cdots>\omega_n>0.
\]
On définit le terme d'ordre par
\[
\mathcal J_{\mathrm{ord}}(\mathbf P)
=
-\sum_{i=0}^n \omega_i\,\mathrm{tr}(P_i\Sigma).
\]
Ce terme impose que les plus grandes directions spectrales soient affectées aux blocs les plus prioritaires.
\end{definition}

\begin{definition}[Objectif architectural conjoint]
On définit l'objectif global par
\[
\mathcal J_{\mathrm{arch}}(\mathbf P,G)
=
\lambda_{\mathrm{leak}}\,\mathcal J_{\mathrm{leak}}(\mathbf P)
+
\lambda_{\mathrm{ord}}\,\mathcal J_{\mathrm{ord}}(\mathbf P)
+
\lambda_{\mathrm{rec}}\,\mathcal E_{\mathrm{rec}}(\mathbf P,G).
\]
En utilisant le théorème précédent, on obtient l'écriture équivalente
\[
\mathcal J_{\mathrm{arch}}(\mathbf P,G)
=
(\lambda_{\mathrm{leak}}+\lambda_{\mathrm{rec}})\,
\mathcal J_{\mathrm{leak}}(\mathbf P)
+
\lambda_{\mathrm{ord}}\,\mathcal J_{\mathrm{ord}}(\mathbf P)
+
\lambda_{\mathrm{rec}}\,\mathcal E_{\mathrm{coll}}(\mathbf P,G).
\]
\end{definition}

\begin{lemma}[Zéro leakage et compatibilité spectrale]
Pour toute partition admissible \(\mathbf P\),
\[
\mathcal J_{\mathrm{leak}}(\mathbf P)\ge 0.
\]
De plus,
\[
\mathcal J_{\mathrm{leak}}(\mathbf P)=0
\quad\Longleftrightarrow\quad
P_i\Sigma=\Sigma P_i
\ \text{ pour tout } i,
\]
c'est-à-dire si et seulement si chaque bloc \(P_i\) est une somme de projecteurs spectraux de \(\Sigma\).
\end{lemma}

\begin{proof}
Par définition, \(\mathcal J_{\mathrm{leak}}\) est une somme de normes au carré, donc elle est non négative.
Si elle vaut \(0\), alors \(P_i\Sigma P_j=0\) pour tous \(i\neq j\), ce qui signifie que \(\Sigma\) est bloc-diagonale dans la décomposition \(\mathbf P\), donc que chaque \(P_i\) commute avec \(\Sigma\).
Réciproquement, si chaque \(P_i\) commute avec \(\Sigma\), alors les blocs hors diagonale sont nuls et \(\mathcal J_{\mathrm{leak}}(\mathbf P)=0\).
\end{proof}

\begin{lemma}[Réécriture télescopique du terme d'ordre]
Soient
\[
m_r=\sum_{i=0}^r d_i
\qquad\text{et}\qquad
Q_r=\sum_{i=0}^r P_i,
\quad r=0,\dots,n,
\]
avec la convention \(Q_{-1}=0\).
Alors
\[
\sum_{i=0}^n \omega_i\,\mathrm{tr}(P_i\Sigma)
=
\omega_n\,\mathrm{tr}(\Sigma)
+
\sum_{r=0}^{n-1}(\omega_r-\omega_{r+1})\,\mathrm{tr}(Q_r\Sigma).
\]
\end{lemma}

\begin{proof}
Comme \(P_i=Q_i-Q_{i-1}\), on a
\[
\sum_{i=0}^n \omega_i\,\mathrm{tr}(P_i\Sigma)
=
\sum_{i=0}^n \omega_i\bigl(\mathrm{tr}(Q_i\Sigma)-\mathrm{tr}(Q_{i-1}\Sigma)\bigr).
\]
Un regroupement télescopique donne
\[
\sum_{i=0}^n \omega_i\,\mathrm{tr}(P_i\Sigma)
=
\omega_n\,\mathrm{tr}(Q_n\Sigma)
+
\sum_{r=0}^{n-1}(\omega_r-\omega_{r+1})\,\mathrm{tr}(Q_r\Sigma).
\]
Comme \(Q_n=\Id\), on obtient la formule annoncée.
\end{proof}

\begin{lemma}[Sélection canonique par l'ordre spectral]
Parmi les partitions admissibles, le terme
\[
\mathcal J_{\mathrm{ord}}(\mathbf P)
=
-\sum_{i=0}^n \omega_i\,\mathrm{tr}(P_i\Sigma)
\]
est strictement minimisé par la partition canonique \(\mathbf P^\star\), et seulement par elle.
\end{lemma}

\begin{proof}
Par le lemme précédent, minimiser \(\mathcal J_{\mathrm{ord}}\) revient à maximiser chaque quantité \(\mathrm{tr}(Q_r\Sigma)\), puisque les coefficients \(\omega_r-\omega_{r+1}\) sont strictement positifs.

Pour tout projecteur orthogonal \(Q\) de rang \(m_r\), le principe de Ky Fan donne
\[
\mathrm{tr}(Q\Sigma)\le \sum_{k=1}^{m_r}\lambda_k,
\]
avec égalité si et seulement si
\[
Q=\sum_{k=1}^{m_r} e_k e_k^\top,
\]
l'unicité étant assurée par la simplicité du spectre.

En l'appliquant à chaque \(Q_r\), on obtient
\[
\mathrm{tr}(Q_r\Sigma)\le \sum_{k=1}^{m_r}\lambda_k,
\]
et l'égalité simultanée pour tout \(r\) impose
\[
Q_r=\sum_{k=1}^{m_r} e_k e_k^\top
\qquad\text{pour tout } r.
\]
Par différence,
\[
P_i=Q_i-Q_{i-1}
=
\sum_{k\in I_i} e_k e_k^\top
=
P_i^\star.
\]
Donc \(\mathbf P^\star\) est l'unique minimiseur de \(\mathcal J_{\mathrm{ord}}\) sur la classe admissible.
\end{proof}

\begin{proposition}[Optimalité du collage inverse-variance]
Sous le modèle de bruit hétéroscédastique, tout estimateur affine non biaisé de la forme
\[
\tilde u_0=\sum_{i=1}^n \beta_i u_0^{(i)},
\qquad
\sum_{i=1}^n \beta_i=1,
\]
vérifie
\[
\mathbb E\|\tilde u_0-u_0^\star\|^2
=
d\sum_{i=1}^n \beta_i^2\tau_i.
\]
Cette quantité est strictement minimisée par
\[
\beta_i^\star
=
\frac{\tau_i^{-1}}{\sum_{j=1}^n\tau_j^{-1}},
\qquad i=1,\dots,n.
\]
\end{proposition}

\begin{proof}
On a
\[
\mathbb E[\tilde u_0-u_0^\star]
=
\sum_{i=1}^n \beta_i\mathbb E[\varepsilon_i]
=
0.
\]
Par indépendance et isotropie,
\[
\mathbb E\|\tilde u_0-u_0^\star\|^2
=
\mathbb E\left\|\sum_{i=1}^n\beta_i\varepsilon_i\right\|^2
=
\sum_{i=1}^n\beta_i^2\,\mathbb E\|\varepsilon_i\|^2
=
d\sum_{i=1}^n\beta_i^2\tau_i.
\]
La fonction
\[
\beta\mapsto d\sum_{i=1}^n\beta_i^2\tau_i
\]
est strictement convexe sur l'hyperplan \(\sum_i\beta_i=1\).  
Le minimiseur unique est donc obtenu par multiplicateurs de Lagrange, ce qui donne
\[
\beta_i^\star
=
\frac{\tau_i^{-1}}{\sum_{j=1}^n\tau_j^{-1}}.
\]
\end{proof}

\begin{theorem}[Optimalité conjointe et unicité du couple partition spectrale + collage pondéré]
Sous les hypothèses précédentes, le problème
\[
\min_{(\mathbf P,G)\ \mathrm{admissible}} \ \mathcal J_{\mathrm{arch}}(\mathbf P,G)
\]
admet un unique minimiseur admissible.
Ce minimiseur est précisément le couple
\[
(\mathbf P^\star,G^\star),
\]
où :

\begin{itemize}[leftmargin=2em]
    \item \(\mathbf P^\star\) est la partition spectrale canonique ;
    \item \(G^\star\) est le collage pondéré inverse-variance, i.e.
    \[
    \bar u_0^\star(w,c)
    =
    \sum_{i=1}^n \alpha_i^\star(w,c)\,u_0^{(i)}(w,c),
    \qquad
    \alpha_i^\star(w,c)
    =
    \frac{\tau_i^{-1}}{\sum_{j=1}^n \tau_j^{-1}}.
    \]
\end{itemize}

En particulier, le couple \((\mathbf P^\star,G^\star)\) est l'unique architecture admissible qui minimise simultanément le leakage inter-modules et l'erreur globale de reconstruction.
\end{theorem}

\begin{proof}
Posons
\[
A=\lambda_{\mathrm{leak}}+\lambda_{\mathrm{rec}}>0.
\]
Alors, pour tout couple admissible \((\mathbf P,G)\),
\[
\mathcal J_{\mathrm{arch}}(\mathbf P,G)
=
A\,\mathcal J_{\mathrm{leak}}(\mathbf P)
+
\lambda_{\mathrm{ord}}\,\mathcal J_{\mathrm{ord}}(\mathbf P)
+
\lambda_{\mathrm{rec}}\,\mathcal E_{\mathrm{coll}}(\mathbf P,G).
\]

\medskip
\noindent
\textbf{Étape 1 : optimisation de la partie partition.}
Comme \(\mathcal J_{\mathrm{leak}}(\mathbf P)\ge 0\) et que \(\mathbf P^\star\) minimise strictement \(\mathcal J_{\mathrm{ord}}(\mathbf P)\) sur la classe admissible, on a pour toute partition admissible
\[
A\,\mathcal J_{\mathrm{leak}}(\mathbf P)+\lambda_{\mathrm{ord}}\,\mathcal J_{\mathrm{ord}}(\mathbf P)
\ge
\lambda_{\mathrm{ord}}\,\mathcal J_{\mathrm{ord}}(\mathbf P)
\ge
\lambda_{\mathrm{ord}}\,\mathcal J_{\mathrm{ord}}(\mathbf P^\star).
\]
De plus,
\[
\mathcal J_{\mathrm{leak}}(\mathbf P^\star)=0
\]
par compatibilité spectrale de la partition canonique.
Ainsi, la partie dépendant de \(\mathbf P\) est minimisée de manière unique par \(\mathbf P^\star\).

\medskip
\noindent
\textbf{Étape 2 : optimisation du collage.}
Une fois \(\mathbf P^\star\) fixé, la partie dépendant de \(G\) se réduit à \(\mathcal E_{\mathrm{coll}}(\mathbf P^\star,G)\).
Par la proposition de réduction du collage,
\[
\argmin_G \mathcal E_{\mathrm{coll}}(\mathbf P^\star,G)
=
\argmin_G \mathbb E\bigl[\|\bar u_0(w,c)-u_0^\star(w,c)\|^2\bigr].
\]
Sous le modèle de bruit hétéroscédastique, la proposition précédente montre que cet ensemble est réduit au seul collage affine non biaisé inverse-variance \(G^\star\), avec
\[
\alpha_i^\star(w,c)=\frac{\tau_i^{-1}}{\sum_{j=1}^n\tau_j^{-1}}.
\]

\medskip
\noindent
\textbf{Étape 3 : unicité globale.}
Si \((\mathbf P,G)\) est un minimiseur global, alors il doit satisfaire simultanément
\[
\mathbf P=\mathbf P^\star
\qquad\text{et}\qquad
G=G^\star,
\]
car l'égalité dans l'étape 1 force \(\mathbf P=\mathbf P^\star\), et l'égalité dans l'étape 2 force \(G=G^\star\).
Réciproquement, \((\mathbf P^\star,G^\star)\) atteint cette borne minimale.
Le minimiseur global est donc unique et vaut exactement \((\mathbf P^\star,G^\star)\).
\end{proof}

\begin{remark}
La décomposition orthogonale précédente est indépendante du mécanisme de collage.
En revanche, l'identification du minimiseur de \(\mathcal E_{\mathrm{coll}}\) avec le collage inverse-variance repose uniquement sur l'hypothèse minimale de compatibilité du mécanisme de collage, qui exprime que \(\mathcal E_{\mathrm{coll}}\) est une transformation strictement croissante de l'erreur de partage.
Le théorème précédent remplace l'hypothèse de décomposition leakage--reconstruction utilisée dans les versions préliminaires du formalisme.
\end{remark}

\begin{corollary}[Optimalité conjointe]
Le formalisme PCH-ESM réalise simultanément :
\begin{enumerate}[leftmargin=2em]
    \item le découplage maximal des modules via la partition spectrale canonique ;
    \item la reconstruction globale de risque minimal via le collage pondéré inverse-variance.
\end{enumerate}
Autrement dit, sous les hypothèses précédentes, le couple \((\mathbf P^\star,G^\star)\) est l'unique minimiseur simultané du leakage inter-modules et de l'erreur de reconstruction dans la classe des architectures admissibles.
\end{corollary}